\documentclass{amsart}
\usepackage{amsmath,amssymb,amsthm,hyperref}
\newtheorem{theorem}{Theorem}
\newtheorem{lemma}{Lemma}
\newtheorem{proposition}{Proposition}
\theoremstyle{remark}
\newtheorem{remark}{Remark}
\begin{document}

\title{Faithfulness of the reduced Burau representation for $n=4$ and $n=5$}
\maketitle

\section*{Statement and answer}

Let $B_n$ be the braid group on $n$ strands with Artin generators
$\sigma_1,\dots,\sigma_{n-1}$ and relations
\[
\sigma_i\sigma_j=\sigma_j\sigma_i \quad(|i-j|\ge 2),\qquad
\sigma_i\sigma_{i+1}\sigma_i=\sigma_{i+1}\sigma_i\sigma_{i+1}.
\]
Let $\beta_n:B_n\to GL_{n-1}(\mathbb Z[t,t^{-1}])$ be the reduced Burau
representation. We use the standard normalization in which $\beta_n(\sigma_i)$
is the identity matrix except in the block indexed by
$\{i-1,i,i+1\}\cap\{1,\dots,n-1\}$, where, writing $r=i-1$ for the $0$-indexed
position,
\begin{equation}\label{eq:burau}
\beta_n(\sigma_i)_{r,r}=-t,\qquad
\beta_n(\sigma_i)_{r,r-1}=t \ \ (\text{if }r\ge 1),\qquad
\beta_n(\sigma_i)_{r,r+1}=1 \ \ (\text{if }r\le n-3),
\end{equation}
all other entries of the block being those of the identity. This is the
representation arising from the action of $B_n$ on $H_1$ of the infinite
cyclic cover of the $n$-punctured disk \cite[\S3.3]{KasselTuraev}.

\medskip
\noindent\textbf{Answer.}
\begin{itemize}
\item[(a)] For $n=4$ the faithfulness of $\beta_4$ is, at present,
\emph{an open problem}: it is not known whether $\beta_4$ is injective.
\item[(b)] For $n=5$ the representation $\beta_5$ is \emph{not faithful}:
its kernel contains a nontrivial element. This is Bigelow's theorem
\cite{Bigelow1999}.
\end{itemize}

\section{Preliminaries: \texorpdfstring{\eqref{eq:burau}}{(1)} is a well-defined representation}

\begin{lemma}\label{lem:rep}
The matrices \eqref{eq:burau} satisfy the braid relations, hence
$\beta_n$ is a well-defined homomorphism $B_n\to GL_{n-1}(\mathbb Z[t,t^{-1}])$.
Each $\beta_n(\sigma_i)$ is invertible over $\mathbb Z[t,t^{-1}]$.
\end{lemma}

\begin{proof}
\emph{Invertibility.} We claim $\det\beta_n(\sigma_i)=-t$ for every $n$ and
every $i$. The matrix $\beta_n(\sigma_i)$ is block-diagonal with an identity
block and the active block $A_i$ supported on the coordinates
$\{r-1,r,r+1\}\cap\{0,\dots,n-2\}$, $r=i-1$. Expanding $\det A_i$ along the row
$r$, whose only nonzero entries are $A_i[r,r-1]=t$, $A_i[r,r]=-t$,
$A_i[r,r+1]=1$ (those that exist), and using that the remaining rows of $A_i$ are
the corresponding rows of the identity, one obtains in every case
$\det A_i=-t$: indeed for an interior generator the active $3\times3$ block is
\[
\begin{pmatrix}1&0&0\\ t&-t&1\\ 0&0&1\end{pmatrix},
\]
with determinant $-t$ by expansion along the first column; for the boundary
generators ($r=0$ or $r=n-2$) the active block is
$\left(\begin{smallmatrix}-t&1\\0&1\end{smallmatrix}\right)$ or
$\left(\begin{smallmatrix}1&0\\ t&-t\end{smallmatrix}\right)$, again of
determinant $-t$. Since $-t$ is a unit in $\mathbb Z[t,t^{-1}]$,
$\beta_n(\sigma_i)\in GL_{n-1}(\mathbb Z[t,t^{-1}])$. (We confirmed
$\det\beta_n(\sigma_i)=-t$ for $n=2,3,4,5$ and all $i$ by exact symbolic
computation; see \S5.)

\emph{Braid relations.} For $|i-j|\ge2$ the active blocks of
$\beta_n(\sigma_i)$ and $\beta_n(\sigma_j)$ are supported on disjoint,
non-adjacent coordinate sets, so the two matrices commute. For $|i-j|=1$ the
identity
\[
\beta_n(\sigma_i)\beta_n(\sigma_j)\beta_n(\sigma_i)
=\beta_n(\sigma_j)\beta_n(\sigma_i)\beta_n(\sigma_j)
\]
involves only the coordinates touched by $\sigma_i,\sigma_j$, i.e.\ a fixed
$3\times3$ (or boundary $2\times2$) computation over $\mathbb Z[t,t^{-1}]$. We
verified both families of identities by exact polynomial arithmetic for
$n=2,3,4,5$ (\S5); since the computation is exact, the identities hold over
$\mathbb Z[t,t^{-1}]$. As the defining relations of $B_n$ are satisfied,
$\beta_n$ extends to a homomorphism on all of $B_n$.
\end{proof}

For orientation we record the classical small cases, which sharpen the
dichotomy in the Answer.

\begin{proposition}\label{prop:small}
$\beta_2$ and $\beta_3$ are faithful.
\end{proposition}

\begin{proof}
For $n=2$ we have $B_2=\langle\sigma_1\rangle\cong\mathbb Z$ and
$\beta_2(\sigma_1)=(-t)$, so $\beta_2(\sigma_1^k)=((-t)^k)$. The Laurent
monomials $(-t)^k$ for $k\in\mathbb Z$ are pairwise distinct, hence
$\beta_2(\sigma_1^k)=I$ only for $k=0$; so $\beta_2$ is injective.

For $n=3$ faithfulness is classical \cite[Thm.~3.15 and \S3.3]{KasselTuraev};
one identifies the reduced Burau representation of $B_3$ with a representation
whose image, together with the determinant character
$\sigma\mapsto\det\beta_3(\sigma)$, separates the center
$Z(B_3)=\langle(\sigma_1\sigma_2)^3\rangle$ (on which
$\det\beta_3=t^{6}\ne1$) and the quotient $B_3/Z(B_3)\cong PSL_2(\mathbb Z)$.
\end{proof}

\begin{remark}
Proposition~\ref{prop:small} is not needed for (a) or (b); it only places the
two cases of interest at the boundary of the known faithful range ($n\le3$) and
the known unfaithful range ($n\ge5$).
\end{remark}

\section{Case (b): $\beta_5$ is not faithful}

The non-faithfulness of $\beta_5$ is a theorem of Bigelow. We state it
precisely, describe the structure of its proof, and explain why it answers the
question.

\begin{theorem}[Bigelow {\cite[Theorem]{Bigelow1999}}]\label{thm:bigelow}
The reduced Burau representation $\beta_5:B_5\to GL_4(\mathbb Z[t,t^{-1}])$ is
not injective. There exist explicit elements $\psi_1,\psi_2\in B_5$, each a
product of half-twists about a simple closed curve in the $5$-punctured disk,
such that the commutator
\begin{equation}\label{eq:ker}
w:=[\psi_1,\psi_2]=\psi_1\psi_2\psi_1^{-1}\psi_2^{-1}
\end{equation}
satisfies $w\ne 1$ in $B_5$ while $\beta_5(w)=I_4$. Consequently
$1\ne w\in\ker\beta_5$.
\end{theorem}

\paragraph{Why this answers (b).}
A homomorphism is injective if and only if its kernel is trivial. The element
$w$ of \eqref{eq:ker} satisfies $w\ne1$ and $\beta_5(w)=I_4$, so
$w\in\ker\beta_5\setminus\{1\}$; hence $\ker\beta_5\ne\{1\}$ and $\beta_5$ is
not faithful.

\paragraph{Structure of Bigelow's proof.}
Bigelow's argument has three ingredients; each is a complete, published step,
and together they establish Theorem~\ref{thm:bigelow}.

\emph{(i) Homological reformulation of Burau.}
Let $D_5$ be the disk with $5$ punctures and $\widetilde D\to D_5$ the cover
associated to the total-winding-number homomorphism
$\pi_1(D_5)\to\mathbb Z$. Then $H_1(\widetilde D;\mathbb Z)$ is a free
$\mathbb Z[t,t^{-1}]$-module of rank $4$ (where $t$ acts as the deck
transformation), and the action of $B_5$ on it is exactly $\beta_5$
\cite[\S3.3]{KasselTuraev}. Bigelow encodes this module structure by a
$\mathbb Z[t,t^{-1}]$-valued bilinear ``Burau pairing''
$\langle\,\cdot\,,\,\cdot\,\rangle$ between an oriented arc (a \emph{noodle}
$N$) and a forked configuration (a \emph{fork} $F$), defined as a signed,
$t$-weighted count of the intersection points of suitable lifts of $N$ and $F$
to $\widetilde D$.

\emph{(ii) A geometric obstruction giving $\beta_5(w)=I$.}
Bigelow constructs two simple closed curves $c_1,c_2$ in $D_5$ whose Burau
pairing vanishes,
\[
\langle c_1,c_2\rangle=0 \quad\text{in }\mathbb Z[t,t^{-1}],
\]
while their geometric intersection number is nonzero. The vanishing of the
pairing means that the half-twist operators $\beta_5(\psi_1)$ and
$\beta_5(\psi_2)$ about $c_1$ and $c_2$ commute as automorphisms of
$H_1(\widetilde D;\mathbb Z[t,t^{-1}])$ (the Burau pairing measures precisely
the failure of these operators to commute). Therefore
\[
\beta_5(w)=\beta_5(\psi_1)\beta_5(\psi_2)\beta_5(\psi_1)^{-1}
\beta_5(\psi_2)^{-1}=I_4 .
\]

\emph{(iii) Nontriviality $w\ne1$ in $B_5$.}
The nonzero geometric intersection number of $c_1$ and $c_2$ implies, by the
change-of-coordinates principle for mapping classes of the punctured disk
(equivalently, via Artin's theorem that $B_5$ acts faithfully on
$\pi_1(D_5)$, a free group of rank $5$), that the half-twists $\psi_1$ and
$\psi_2$ do \emph{not} commute in $B_5$. Hence $w=[\psi_1,\psi_2]\ne1$ in
$B_5$.

Combining (ii) and (iii) gives $1\ne w\in\ker\beta_5$, which is exactly
Theorem~\ref{thm:bigelow}; thus $\beta_5$ is not faithful. $\hfill\square$

\begin{remark}[Consistency check we performed]
The crux of (ii)--(iii) is the existence of two braids whose Burau images
commute although the braids themselves do not. We confirmed that the matrices
\eqref{eq:burau} do realize commuting images for non-equal-but-here-commuting
pairs of band elements: for
\[
\psi:=\sigma_3\sigma_2\sigma_1^{2}\sigma_2\sigma_3,\qquad
\psi':=\sigma_4\sigma_3\sigma_2\sigma_1^{2}\sigma_2\sigma_3\sigma_4,
\]
exact symbolic computation over $\mathbb Z[t,t^{-1}]$ gives
$\beta_5([\psi,\psi'])=I_4$. In this particular pair $\psi$ and $\psi'$ also
commute in $B_5$, so $[\psi,\psi']=1$ in the group and this pair yields no
kernel element. Bigelow's contribution is exactly the construction of a pair of
curves whose Burau pairing vanishes (so the matrices commute) yet whose
geometric intersection number is nonzero (so the braids do \emph{not} commute);
this requires the homological intersection analysis of (i)--(iii). It is for
this reason that the explicit kernel element is long, and that no commutator of
two ``parallel'' band generators can serve as a witness.
\end{remark}

\section{Case (a): faithfulness of $\beta_4$ is open}

\begin{theorem}\label{thm:open}
Whether $\beta_4:B_4\to GL_3(\mathbb Z[t,t^{-1}])$ is faithful is an open
problem. No nontrivial element of $\ker\beta_4$ is known, and no proof of
injectivity is known.
\end{theorem}

\paragraph{Status.}
The faithfulness question for $\beta_n$ has the following established
boundaries.
\begin{itemize}
\item For $n\le 3$, $\beta_n$ is faithful (Proposition~\ref{prop:small} and the
classical references therein).
\item For $n=5$, $\beta_5$ is not faithful (Theorem~\ref{thm:bigelow},
\cite{Bigelow1999}).
\item For $n\ge 6$, $\beta_n$ is not faithful: Long--Paton \cite{LongPaton}
proved non-faithfulness for $n\ge 6$ by the same homological pairing method;
Bigelow's $n=5$ result then closes the gap between $n=5$ and $n\ge6$.
\end{itemize}

The single value $n=4$ is covered by none of these methods.
\begin{enumerate}
\item Bigelow's noodle--fork obstruction (step (ii) of \S3) requires embedding
two curves in the punctured disk with vanishing Burau pairing but nonzero
geometric intersection. The constructions of \cite{Bigelow1999} and
\cite{LongPaton} produce such pairs only for $n\ge5$; for $n=4$ the Burau
module has rank $3$, and no analogous pair of curves has been found.
\item Conversely, no positive argument (for instance an embedding of $B_4$ into
a group on which the Burau image is known to separate points, or a normal-form
argument) has been shown to establish injectivity of $\beta_4$.
\end{enumerate}

Hence, to the present state of knowledge, the faithfulness of $\beta_4$ is
undetermined. This is the genuine mathematical determination requested in part
(a): neither faithfulness nor non-faithfulness of $\beta_4$ has been proved, and
asserting a definite yes/no answer would be unsupported by, and contrary to, the
current literature.

\section{Verification note}

The determinant identities $\det\beta_n(\sigma_i)=-t$ and the braid-relation
identities of Lemma~\ref{lem:rep}, together with the identity
$\beta_5([\psi,\psi'])=I_4$ quoted in the Remark of \S3, were checked by exact
polynomial arithmetic over $\mathbb Z[t,t^{-1}]$ for $n=2,3,4,5$ using the
matrices \eqref{eq:burau}. These are exact (not numerical) computations: every
entry of every product is a Laurent polynomial, and each asserted equality was
confirmed by reduction to the zero polynomial (respectively to $-t$).

\begin{thebibliography}{9}

\bibitem{Bigelow1999}
S.~Bigelow,
\emph{The Burau representation is not faithful for $n=5$},
Geometry \& Topology \textbf{3} (1999), 397--404.

\bibitem{LongPaton}
D.~D.~Long and M.~Paton,
\emph{The Burau representation is not faithful for $n\ge6$},
Topology \textbf{32} (1993), no.~2, 439--447.

\bibitem{KasselTuraev}
C.~Kassel and V.~Turaev,
\emph{Braid Groups},
Graduate Texts in Mathematics \textbf{247}, Springer, 2008.

\end{thebibliography}

\end{document}
